\section{Finite obstructions, compatible lifts, and deformation}
\label{sec:lifting}

The arithmetic results above concern particular families. We now study the
underlying correction problem, from finite polynomial lifts to continuous
parameter charts and quantitative deformation. Throughout, we distinguish
three objects: a residual obstruction, the set of correcting digits, and a
root in a completion. A vanishing obstruction at one precision gives a lift
to the next. Continuing indefinitely requires additional hypotheses.

Unless a different coefficient ring is specified, a polynomial system is a
finite list of integer polynomials. All divisibility conditions on vectors are
coordinatewise. For a positive integer $q$, write $S_q=\Z/q\Z$ and use an
overline for reduction modulo $q$. For a prime $p$, write
$\F_p=S_p$, normalize $|p|_p=p^{-1}$, and put
\[
 B_L(x)=x+p^L\Zp^d.
\]
An integer center is included in this notation by its natural image in
$\Zp^d$. All dimensions below are finite and may be zero. A statement over
$S_q$ allows composite $q$ and $q=1$ when explicitly stated. A cardinality
formula in terms of vector-space rank uses prime $p$.

\subsection{The polynomial correction equation}

For $F=(F_1,\ldots,F_e)\in\Z[X_1,\ldots,X_d]^e$, define its evaluated
Jacobian by the formal partial derivatives:
\[
 DF(x)_{ij}=(\partial_jF_i)(x),\qquad
 J_x=DF(x)\bmod q:S_q^d\longrightarrow S_q^e.
\]
The obstruction therefore uses the Jacobian of the given polynomials. The
following algebraic Taylor identity relates this linear map to polynomial
evaluation.

\begin{lemma}[Polynomial Taylor remainder]\label{lem:polynomial-taylor}
Let $R$ be a commutative ring, $P\in R[X_1,\ldots,X_d]$, $a\in R^d$, and
$c\in R$. There is a polynomial $Q\in R[U_1,\ldots,U_d]$, with $Q(0)=0$,
such that for every $u\in R^d$,
\begin{equation}\label{eq:polynomial-taylor}
 P(a+cu)=P(a)+c\sum_{j=1}^d(\partial_jP)(a)u_j+c^2Q(u).
\end{equation}
In particular, when $R=\Z$ and $a,u,c$ are integral, subtracting the first
two terms on the right from $P(a+cu)$ gives a multiple of $c^2$.
\end{lemma}

\begin{proof}
Apply polynomial induction using constants, addition, and multiplication by a
variable. Constants have zero remainder, and remainders add. For the
multiplication step, suppose the assertion is known for $P$ and write
$D_P(U)=\sum_j(\partial_jP)(a)U_j$. Multiplication of its expansion by
$a_i+cU_i$ gives constant term $P(a)a_i$, linear term
$c(a_iD_P(U)+P(a)U_i)$, and remainder
\[
 c^2\bigl(Q(U)(a_i+cU_i)+D_P(U)U_i\bigr).
\]
The product rule identifies the linear term with the derivative of
$PX_i$. The displayed remainder vanishes at $U=0$. This constructs the
polynomial remainder before evaluation at $u$, even over rings where
division by $c$ would be invalid.
\end{proof}

The general polynomial identity is formalized by
\lean{Komorebi.PolynomialTaylorRing.exists_eval_scaled_remainder}.
Its integer divisibility form is
\lean{Komorebi.PolynomialTaylor.taylor_dvd}.

\begin{theorem}[Finite polynomial obstruction]\label{thm:finite-polynomial-obstruction}
Let $q\ge1$, $L\ge1$, $x\in\Z^d$, and suppose
$F(x)\in q^L\Z^e$. Define the exact integer residual $r_x$ by
$F(x)=q^Lr_x$, and put
\[
 \mathcal O_x=S_q^e/\im J_x,\qquad o_x=[\overline{r_x}].
\]
Then $\mathcal O_x$ is finite and
\begin{equation}\label{eq:finite-obstruction}
 o_x=0\quad\Longleftrightarrow\quad
 \exists y\in\Z^d:\quad
 y\equiv x\pmod{q^L},\qquad F(y)\equiv0\pmod{q^{L+1}}.
\end{equation}
No surjectivity or nonsingularity assumption is required.
\end{theorem}

\begin{proof}
Every point in the lower residue class has the form $y=x+q^Lt$ with
$t\in\Z^d$. By Lemma~\ref{lem:polynomial-taylor} and $2L\ge L+1$,
\[
 F(x+q^Lt)\equiv q^L\bigl(r_x+DF(x)t\bigr)\pmod{q^{L+1}}.
\]
The integer $q^L$ is nonzero. Cancelling it in the corresponding divisibility
condition proves
\begin{equation}\label{eq:digit-correction}
 F(x+q^Lt)\equiv0\pmod{q^{L+1}}
 \quad\Longleftrightarrow\quad
 \overline{r_x}+J_x\overline t=0.
\end{equation}
Integer vectors reduce onto $S_q^d$. The equation on the right is solvable
exactly when $\overline{r_x}$ belongs to $\im J_x$, since that image is
closed under negation. This is precisely $o_x=0$. Finiteness follows because
$\mathcal O_x$ is a quotient of the finite module $S_q^e$.
\end{proof}

The formal theorem
\lean{Komorebi.PolynomialObstruction.finite_polynomial_obstruction}
presents the image using the integer-domain map
$\Z^d\xrightarrow{DF(x)}\Z^e\to S_q^e$. Reduction onto $S_q^d$ identifies
that image with $\im J_x$. The class-preserving comparison is explicit in
\lean{Komorebi.PolynomialLiftCounts.obstructionEquiv}; the change of
presentation does not change the distinguished residual class.

\begin{proposition}[Affine depth zero]\label{prop:affine-obstruction}
Let $F(x)=Ax-b$ for an integer linear map $A:\Z^d\to\Z^e$. The criterion
of Theorem~\ref{thm:finite-polynomial-obstruction} holds for every $L\ge0$,
with $J=A\bmod q$.
\end{proposition}

\begin{proof}
The affine identity
$F(x+q^Lt)=F(x)+q^LAt$ has no quadratic remainder. The proof of
\eqref{eq:digit-correction} therefore needs no inequality involving $2L$.
The residual and the finite quotient are obtained as above.
\end{proof}

The affine statement is formalized by
\lean{Komorebi.NativeAffine.finite_affine_obstruction}.
Positive depth is essential in the nonlinear theorem: for $F(X)=X^2-1$,
$q=2$, $L=0$, and $x=0$, the Jacobian is zero and the residual class is
nonzero, yet $y=1$ is a root modulo two in the same residue modulo one.

\begin{proposition}[Independence of the integer representative]
\label{prop:representative-independence}
Under the hypotheses of Theorem~\ref{thm:finite-polynomial-obstruction},
let $x'=x+q^Ls$ and write $F(x')=q^Lr_{x'}$. Then
\[
 J_{x'}=J_x,\qquad
 \overline{r_{x'}}=\overline{r_x}+J_x\overline s.
\]
Consequently the two residual classes coincide in the common quotient
$S_q^e/\im J_x$.
\end{proposition}

\begin{proof}
Polynomial evaluation respects coordinatewise congruence. Apply this first
to every formal partial derivative: $x'\equiv x\pmod q$ gives
$J_{x'}=J_x$. Next apply \eqref{eq:polynomial-taylor} to $F$ at scale $q^L$.
Subtract the two exact residual identities and cancel $q^L$ to obtain the
displayed relation modulo $q$. Its additional term belongs to the common
Jacobian image, so it vanishes in the quotient.
\end{proof}

This comparison is formalized by
\lean{Komorebi.PolynomialRepresentative.representative_independence}.
This proves independence of the integer representative within a fixed
polynomial presentation. It makes no identification of the raw cokernels
of different equation lists. For example, adjoining a zero equation leaves
the roots unchanged while adding an $S_q$ summand to the raw cokernel.

\subsection{One correction shared by two systems}

The next obstruction concerns compatibility between individually attainable
corrections. Its algebraic form works over an arbitrary commutative ring.
Let $U,V,W$ be $R$-modules and let $J:U\to V$, $K:U\to W$ be linear.
Define
\begin{equation}\label{eq:common-quotient}
 C(J,K)=U/(\ker J+\ker K).
\end{equation}
An attainable value $Ju$ has label $[u]\in C(J,K)$; this is independent of
the choice of its preimage because $\ker J$ is killed in the quotient.
Similarly, $Kv$ has label $[v]$. Denote these maps by
$\lambda_J:\im J\to C(J,K)$ and $\lambda_K:\im K\to C(J,K)$.

\begin{theorem}[Common correction and fiber product]\label{thm:common-correction}
For every $u,v\in U$,
\begin{equation}\label{eq:common-correction}
 [u]=[v]\text{ in }C(J,K)
 \quad\Longleftrightarrow\quad
 \exists w\in U:\ Jw=Ju,\quad Kw=Kv.
\end{equation}
Both label maps are surjective, and the image of the joint map is exactly
the fiber product:
\begin{equation}\label{eq:joint-fiber-product}
 \im\bigl(U\to\im J\times\im K,\ u\mapsto(Ju,Ku)\bigr)
 =\im J\mathbin{\times}_{C(J,K)}\im K.
\end{equation}
If either $V$ or $W$ is finite, then $C(J,K)$ is finite.
\end{theorem}

\begin{proof}
If $[u]=[v]$, write $u-v=s+t$ with $s\in\ker J$ and $t\in\ker K$.
Then $w=u-s=v+t$ satisfies both required equalities. Conversely, a common
$w$ gives
\[
 u-v=(u-w)+(w-v)\in\ker J+\ker K.
\]
Every class $[u]$ is the label of $Ju$ and of $Ku$, proving surjectivity.
For attainable $a=Ju$ and $b=Kv$, the first equivalence says exactly that
equal labels are equivalent to a simultaneous preimage. This proves
\eqref{eq:joint-fiber-product}. Finally, a finite target has finite image,
and the associated surjective label map makes $C(J,K)$ finite.
\end{proof}

These statements are
\lean{Komorebi.SharedCorrection.quotient_mk_eq_iff_common} and
\lean{Komorebi.SharedCorrection.joint_range_eq_fiber}, together with the
finite-target corollaries in the same module. The quotient also has the
expected universal property: every linear map from $U$ annihilating both
kernels factors uniquely through $C(J,K)$. Equivalently, compatible linear
observations on $\im J$ and $\im K$ factor uniquely through their label
maps. This is formalized by
\lean{Komorebi.SharedCorrection.common_paired_universal}.

\begin{corollary}[Simultaneous polynomial lifting]\label{cor:shared-polynomial}
Let $F\in\Z[X_1,\ldots,X_d]^e$ and
$G\in\Z[X_1,\ldots,X_d]^f$ have a common lower root $x$ modulo $q^L$,
where $q\ge1$ and $L\ge1$. Use the integer-domain reduced Jacobians
$J:\Z^d\to S_q^e$ and $K:\Z^d\to S_q^f$. Suppose $u,v\in\Z^d$ are
separately valid corrections for $F$ and $G$, respectively. Then
\[
 [u]=[v]\text{ in }\Z^d/(\ker J+\ker K)
\]
if and only if a single $y\equiv x\pmod{q^L}$ solves both systems modulo
$q^{L+1}$. Each separate label is independent of the chosen valid
correction. The common quotient is finite.
\end{corollary}

\begin{proof}
For each system, \eqref{eq:digit-correction} translates a valid next root
into a linear correction equation. Once $u$ corrects $F$, another digit
$w$ corrects $F$ precisely when $Jw=Ju$. The analogous condition for $G$
is $Kw=Kv$. Theorem~\ref{thm:common-correction} now gives the criterion.
Two valid corrections of one block differ by a vector in its Jacobian
kernel, so they have equal labels. Finiteness follows from either finite
residue target. An unconditional formulation quantifies over separately
valid $u,v$ with equal labels, thereby including their existence in the
criterion rather than assuming it.
\end{proof}

See \lean{Komorebi.PolynomialShared.joint_lift_iff_class_eq} and
\lean{Komorebi.PolynomialShared.finite_polynomial_pair_compatibility}.

For affine integer systems $F(X)=AX-b$ and $G(X)=BX-c$, where
$A:\Z^d\to\Z^e$ and $B:\Z^d\to\Z^f$ are linear,
Corollary~\ref{cor:shared-polynomial} holds for $q\ge1$ and every
$L\ge0$, including depth zero. Here $x$ is a common root modulo $q^L$,
and $J$ and $K$ are the reductions of $A$ and $B$ with domain $\Z^d$.
Given separately valid corrections $u,v$, a simultaneous next lift exists
if and only if their classes in $\Z^d/(\ker J+\ker K)$ agree. The quotient
is finite, and each label is independent of the chosen valid correction.
Indeed, Proposition~\ref{prop:affine-obstruction} supplies the correction
equation without a positive-depth restriction, and the proof of
Theorem~\ref{thm:common-correction} then applies unchanged. The formal
statement is \lean{Komorebi.AffinePair.joint_lift_iff_class_eq}.
Quantifying over separately valid corrections with equal labels also
includes their existence in the criterion, as in
\lean{Komorebi.AffinePair.finite_affine_pair_compatibility}.
Separate liftability alone is insufficient. At $x=0$ modulo two,
$F(X)=X$ and $G(X)=X-2$ both lift to modulus four, but their correcting
digits are respectively $0$ and $1$ modulo two. No single residue solves
both equations modulo four.

\subsection{Correction torsors and exact residue counts}

Integer representatives and residue classes at the next level must also be
distinguished. For $q\ge1$ and $L\ge0$, the map
\[
 S_q^d\longrightarrow
 \{y\in(\Z/q^{L+1}\Z)^d:y\equiv x\pmod{q^L}\},
 \qquad t\longmapsto x+q^L\widetilde t
\]
is a bijection, where $\widetilde t$ is any integer representative.
Independence of the representative and injectivity both follow by cancelling
$q^L$ in integer divisibility; surjectivity follows by expressing the lower
congruence as a multiple of $q^L$. This is the equivalence
\lean{Komorebi.ResidueLiftDigits.digitEquivResidue}.

At positive polynomial depth, restrict it to the correction fiber
\[
 \mathcal T_x=\{t\in S_q^d:\overline{r_x}+J_xt=0\}
\]
and the lift set
\[
 \mathcal L_{x,L}=
 \{y\in(\Z/q^{L+1}\Z)^d:y\equiv x\pmod{q^L},\ F(y)=0\}.
\]
The polynomial predicate in the second display is evaluation in the residue
ring. Evaluation on canonical integer representatives gives the same
predicate, because integer polynomial evaluation respects congruence.

\begin{theorem}[Torsor and lift cardinalities]\label{thm:lift-torsor-count}
Under the hypotheses of Theorem~\ref{thm:finite-polynomial-obstruction},
there is an equivalence $\mathcal T_x\simeq\mathcal L_{x,L}$.
Every nonempty fiber is a torsor for $\ker J_x$ under addition of
correction digits. If $q=p$ is prime and
$\rho_x=\rank_{\F_p}J_x$, then
\begin{equation}\label{eq:lift-count}
 \#\mathcal L_{x,L}=
 \begin{cases}p^{d-\rho_x},&o_x=0,\\0,&o_x\ne0,\end{cases}
 \qquad
 \#\mathcal O_x=p^{e-\rho_x}.
\end{equation}
\end{theorem}

\begin{proof}
Equation~\eqref{eq:digit-correction} identifies exactly which points of the
residue fiber satisfy the equations, proving the restricted equivalence.
If $u,t\in\mathcal T_x$, subtraction gives $J_x(t-u)=0$. Conversely,
adding any element of $\ker J_x$ preserves the correction equation. The
unique kernel vector carrying $u$ to $t$ is $t-u$, so the action is free
and transitive. Choosing one valid digit identifies the fiber with the
kernel, without making that choice part of its translation action.

For prime $p$, rank-nullity gives
$\dim\ker J_x=d-\rho_x$ and
$\dim(\F_p^e/\im J_x)=e-\rho_x$. The cardinality of an
$n$-dimensional $\F_p$-space is $p^n$. Combine these identities with the
obstruction criterion and the residue equivalence.
\end{proof}

The cardinality formula is formalized by
\lean{Komorebi.PolynomialLiftCounts.finite_polynomial_lift_count}.
The explicit torsor and chosen-point equivalence are developed in
\lean{Komorebi.FiniteFiberCounts.translate_exists_unique} and
\lean{Komorebi.PolynomialLiftCounts.liftEquivKernel}.
For composite $q$, the torsor and kernel cardinality remain valid, while the
vector-space rank formula is unavailable. Multiplication by two on
$\Z/4\Z$ already has a kernel of size two.

\subsection{Smooth compatible towers and exact roots}

Surjectivity removes every residual obstruction and persists throughout a
positive-depth residue ball. This supplies the additional hypothesis needed
to iterate the finite theorem.

\begin{theorem}[Compatible polynomial tower]\label{thm:smooth-tower}
Let $q\ge1$, $L\ge1$, $F\in\Z[X_1,\ldots,X_d]^e$, and
$x\in\Z^d$ satisfy $F(x)\in q^L\Z^e$. Suppose $J_x:S_q^d\to S_q^e$
is surjective. There is a single sequence $z_n\in\Z^d$ such that
\begin{align*}
 z_0&=x,& F(z_n)&\in q^{L+n}\Z^e,\\
 z_{n+1}-z_n&\in q^{L+n}\Z^d,& J_{z_n}&=J_x.
\end{align*}
For $a\le b$, it also satisfies $z_b-z_a\in q^{L+a}\Z^d$.
\end{theorem}

\begin{proof}
At level $N=L+n$, divide the residual by $q^N$. Surjectivity supplies
$t\in S_q^d$ with $J_{z_n}t=-\overline{F(z_n)/q^N}$. Choose an integer
representative and apply \eqref{eq:digit-correction} to obtain $z_{n+1}$.
The new point is congruent to $z_n$ modulo $q$. Evaluation of each formal
partial derivative therefore gives $J_{z_{n+1}}=J_{z_n}$.

Construct the sequence recursively, retaining the root condition and the
fixed Jacobian at each step. Surjectivity at every later stage then follows
from the single initial assumption. The cross-level congruence follows by
induction, adding consecutive differences whose divisibilities all imply divisibility
by $q^{L+a}$.
\end{proof}

The full tower statement is
\lean{Komorebi.PolynomialIntegerTower.exists_compatible_tower_strong}.
The proof constructs a compatible path using classical choice. It does not
claim an optimized procedure for selecting its digits.

\begin{theorem}[Completion and smooth branching]\label{thm:smooth-padic-root}
In Theorem~\ref{thm:smooth-tower}, let $q=p$ be prime. The tower can be
chosen with a limit $a\in\Zp^d$ such that
\[
 F(a)=0,\qquad z_n\longrightarrow a,\qquad
 a-z_n\in p^{L+n}\Zp^d\quad(n\ge0).
\]
At each stage, the number of distinct next-level residue roots over $z_n$
is exactly $p^{d-e}$, and the finite obstruction module has one element.
If $d=e$, there is exactly one root of $F$ in $B_L(x)$.
\end{theorem}

\begin{proof}
Cross-level compatibility bounds the supremum distance between later terms
by $p^{-(L+n)}$. Completeness gives a limit, and closed norm bounds retain
the congruence to every $z_n$. View the coefficients of the original integer
polynomials in $\Zp$. Their evaluations satisfy
both $F_i(z_n)\to F_i(a)$, by continuity, and $F_i(z_n)\to0$, by their
growing divisibility. Uniqueness of limits gives exact vanishing.

Surjectivity gives $\rank J_x=e\le d$. Theorem~\ref{thm:lift-torsor-count}
then gives $p^{d-e}$ next residue roots and obstruction cardinality one at
every stage.

For square uniqueness, let $a,b\in B_L(x)$ be arbitrary exact roots.
Their canonical integer truncations are genuine roots in the respective
finite residue rings, since evaluation commutes with reduction. They agree
at depth $L$. If they agree at depth $N\ge L$, both next residues belong
to the same lift set, which has cardinality $p^{d-d}=1$. Induction gives
agreement at every depth. The residue maps separate $p$-adic integers, so
$a=b$. Existence was already supplied by the tower.
\end{proof}

The endpoint construction is
\lean{Komorebi.PolynomialPadicLift.exists_padic_root_with_tower}, the
branching statement is
\lean{Komorebi.PolynomialSmoothCounts.exists_regular_tower}, and uniqueness
for all roots in the ball is
\lean{Komorebi.PolynomialPadicUnique.exists_unique_padic_root}.

\begin{remark}\label{rem:singular-tower-boundary}
The one-step singular criterion and the smooth infinite theorem have
different hypotheses. The polynomial $X^2-5$ has odd roots modulo two and
four but has no root modulo eight. Thus a previous zero obstruction does
not force the next one to vanish. Also, rectangular surjectivity alone
does not imply uniqueness: $F(X,Y)=X$ has the distinct dyadic roots
$(0,0)$ and $(0,2)$ in the same residue ball modulo two. These examples show
why the hypotheses of the tower theorem and its uniqueness clause are needed.
\end{remark}

\subsection{Coefficient stability and continuous root charts}

\begin{proposition}[Smooth coefficient stability]\label{prop:smooth-coefficient-stability}
Let $p$ be prime, $M\ge1$, $F,H\in\Z[X_1,\ldots,X_d]^e$, and
$c\in p^M\Z$. Suppose $DF(x)\bmod p$ is surjective for an integer
$x\in\Z^d$. If $a\in\Zp^d$ satisfies $F(a)=0$ and
$a\equiv x\pmod p$, then there is a root $b$ of $G=F+cH$ with
$b-a\in p^M\Zp^d$.

If $d=e$ and $M\ge L\ge1$, every pair of roots $a,b\in B_L(x)$ of
$F$ and $G$, respectively, satisfies $b-a\in p^M\Zp^d$. The perturbed
root is unique in that ball whenever it exists.
\end{proposition}

\begin{proof}
Since $p\mid c$, differentiation gives
$DG(x)\equiv DF(x)\pmod p$. Truncate $a$ to an integer vector $y$ at
depth $M$. Then $F(y)\equiv0\pmod{p^M}$, hence also
$G(y)\equiv0\pmod{p^M}$. Congruence of $y$ and $x$ modulo $p$ preserves
the reduced Jacobian. The smooth root theorem gives
$b\equiv y\equiv a\pmod{p^M}$.

In the square case this constructed root lies in $B_L(x)$ because $M\ge L$.
Square uniqueness for $G$ identifies it with every competing root in that
ball. The construction therefore gives the precision assertion for an
arbitrary root there.
\end{proof}

See \lean{Komorebi.PolynomialRootStability.exists_nearby_perturbed_root}
and \lean{Komorebi.PolynomialRootStability.perturbed_roots_congruent}.
The scalar $c$ may be negative or zero. The weighted theorems below allow
arbitrary integral $p$-adic coefficients and give an explicit first-order
displacement.

\begin{theorem}[Continuous charts with a selected Jacobian minor]
\label{thm:continuous-parameter-chart}
Let $p$ be prime, $L\ge1$,
$F\in\Z[X_1,\ldots,X_{e+k}]^e$, and $x\in\Z^{e+k}$ satisfy
$F(x)\in p^L\Z^e$. Suppose the full reduced Jacobian is surjective.
There is a permutation $\sigma$ of the coordinates such that, writing
$\sigma x=(y_0,t_0)$ with $e$ solved coordinates and $k$ parameter
coordinates, the parameter ball
$T=t_0+p^L\Zp^k$ is homeomorphic to the entire root locus
\[
 \mathcal R=\{z\in B_L(x):F(z)=0\}.
\]
The inverse homeomorphism is the actual free-coordinate projection
$z\mapsto\operatorname{free}(\sigma z)$. If $h:T\to\mathcal R$ is the
forward map, then for every $N\ge0$,
\begin{equation}\label{eq:chart-congruence}
 h(t)\equiv h(u)\pmod{p^N}
 \quad\Longleftrightarrow\quad t\equiv u\pmod{p^N}.
\end{equation}
Equivalently, $\sigma h(t)=(\phi(t),t)$ for a continuous
$\phi:T\to y_0+p^L\Zp^e$ whose graph is the entire reordered root locus.
\end{theorem}

\begin{proof}
First consider a square integer polynomial map $P$ whose reduced Jacobian
at $x$ is surjective. Its Jacobian is invertible over $\F_p$. A digit
comparison using Taylor expansion shows that equal outputs modulo
$p^{N+1}$ force equal inputs modulo $p^{N+1}$ whenever the inputs already
agree modulo $p^N$ and belong to $B_1(x)$. Induction and integer truncation
give, for arbitrary $a,b\in B_1(x)$,
\[
 P(a)\equiv P(b)\pmod{p^N}
 \quad\Longleftrightarrow\quad a\equiv b\pmod{p^N}.
\]

More generally, for a rectangular map with surjective reduced Jacobian,
\begin{equation}\label{eq:polynomial-ball-image}
 P(B_L(x))=P(x)+p^L\Zp^e.
\end{equation}
One inclusion follows from polynomial congruence. For integer targets in
the right-hand ball, subtract the target from the constant terms and use
the smooth root theorem. For arbitrary $p$-adic targets, approximate by
integer truncations. The input ball is compact, so its continuous image is
closed. It contains all these approximants and hence their limit. This
argument does not assume that independently chosen preimages form a
compatible sequence.

For square $P$, the resulting bijection between the two compact residue
balls is continuous, and its target is Hausdorff. It is therefore a
homeomorphism, with the same congruence preservation for its inverse.

Now suppose a chosen set of $e$ columns of $DF(x)\bmod p$ is independent,
and put those coordinates first. Apply the square result to the augmented
polynomial map
\[
 \Psi(Y,T)=(F(Y,T),T).
\]
Its reduced Jacobian has upper blocks $(D_YF,D_TF)$ and lower blocks
$(0,I_k)$, so it is invertible. For $t\in T$, the target $(0,t)$ belongs
to $\Psi(x)+p^L\Zp^{e+k}$ because $F(x)\equiv0\pmod{p^L}$.
The inverse of $\Psi$ at that target gives the unique root with parameters
$t$. The inverse congruence property gives one direction of
\eqref{eq:chart-congruence}; projection onto the parameter coordinates
gives the other.

Finally, surjectivity of the original rectangular Jacobian guarantees a
basis among its columns. Extend the corresponding injection of coordinate
indices to a permutation $\sigma$. Renaming the variables by
$\sigma^{-1}$ transports polynomial evaluation and every formal partial
derivative. The fixed-coordinate construction therefore applies to the
original system after undoing the permutation.
\end{proof}

The main theorem is
\lean{Komorebi.PolynomialGeneralChart.exists_parameter_homeomorph}.
The square local inverse and surjectivity for arbitrary targets are
\lean{Komorebi.PolynomialLocalInverse.localHomeomorph} and
\lean{Komorebi.PolynomialLocalSurjective.exists_padic_preimage}.
The center $x$ is only an approximate root: the chart need not satisfy
$\phi(t_0)=y_0$ unless $F(x)=0$ exactly. The assertion selects suitable
coordinates; it does not permit an arbitrary prescribed solved block.

\subsection{Prescribed parameters and their obstruction}

Suppose now that the split into solved variables and parameter variables
is fixed in advance. Let
\[
 F\in\Z[Y_1,\ldots,Y_m,T_1,\ldots,T_k]^e,\qquad x=(y_0,t_0),
\]
and assume $F(x)=q^Lr_x$ with $q\ge1$ and $L\ge1$. Write the reduced
Jacobian as
\[
 DF(x)\bmod q=(J_x\ K_x),\qquad
 J_x:S_q^m\to S_q^e,\quad K_x:S_q^k\to S_q^e.
\]
For a prescribed parameter digit $v\in S_q^k$, let $\mathcal L_{x,L,v}$
be the set of next residue roots over $x$ whose parameter coordinates are
$t_0+q^L\widetilde v$ modulo $q^{L+1}$.

\begin{theorem}[Prescribed-parameter obstruction and count]
\label{thm:prescribed-obstruction}
Define
\[
 o_x(v)=[\overline{r_x}+K_xv]\in S_q^e/\im J_x.
\]
Then $\mathcal L_{x,L,v}$ is nonempty if and only if $o_x(v)=0$.
When nonempty, it is a torsor for $\ker J_x$. If $q=p$ is prime and
$\rho_x=\rank_{\F_p}J_x$, then
\begin{equation}\label{eq:prescribed-count}
 \#\mathcal L_{x,L,v}
 =\begin{cases}p^{m-\rho_x},&o_x(v)=0,\\0,&o_x(v)\ne0,\end{cases}
 \qquad \#(\F_p^e/\im J_x)=p^{e-\rho_x}.
\end{equation}
The criterion holds uniformly for every prescribed $v$ without a
surjectivity hypothesis.
\end{theorem}

\begin{proof}
A candidate for the next lift has the form
$x+q^L(\widetilde u,\widetilde v)$, with solved digit $u\in S_q^m$.
Its Taylor correction equation is
\[
 \overline{r_x}+J_xu+K_xv=0.
\]
Solvability is exactly vanishing of the displayed cokernel class. Restrict
the residue-digit equivalence to the prescribed parameter coordinates.
This identifies the lift fiber with the solved-digit fiber.
Translation by $\ker J_x$ is free and transitive on any nonempty such
fiber. The prime-field count follows from rank-nullity as before.
\end{proof}

The formal statements start from the lower root and construct its residual:
\lean{Komorebi.PolynomialPrescribedObstruction.finite_prescribed_obstruction}
and \lean{Komorebi.PolynomialPrescribedObstruction.finite_prescribed_lift_count}.
For $F(Y,T)=T$ at the dyadic origin, the full Jacobian $(0\ 1)$ is
surjective but $J_x=0$. Prescribing $v=1$ forces $T=2\pmod4$, so no root
exists; prescribing $v=0$ gives two choices of $Y$ modulo four over the
initial zero residue. This example explains why the criterion depends on
the solved block, rather than merely on the full Jacobian.

\begin{theorem}[Every compatible integer parameter path]
\label{thm:prescribed-integer-tower}
In the preceding notation, assume $J_x$ is surjective. For every supplied
integer path $t_n\in\Z^k$ satisfying
\[
 t_0=\operatorname{free}(x),\qquad
 t_{n+1}-t_n\in q^{L+n}\Z^k,
\]
there is a compatible polynomial tower $z_n\in\Z^{m+k}$ with
\[
 z_0=x,\quad \operatorname{free}(z_n)=t_n,\quad
 F(z_n)\in q^{L+n}\Z^e,\quad
 z_{n+1}-z_n\in q^{L+n}\Z^{m+k}.
\]
The free-coordinate equality is exact in integers. Both reduced Jacobian
blocks remain fixed, and $z_b-z_a\in q^{L+a}\Z^{m+k}$ whenever $a\le b$.
\end{theorem}

\begin{proof}
At level $N=L+n$, write the supplied parameter increment exactly as
$t_{n+1}-t_n=q^Nw$ with $w\in\Z^k$. Surjectivity of the solved block
gives a digit $u$ satisfying
$J_{z_n}u=-\overline{F(z_n)/q^N}-K_{z_n}\overline w$.
Choose an integer representative $\widetilde u$ and set
$z_{n+1}=z_n+q^N(\widetilde u,w)$. Its parameter coordinates equal the
supplied target exactly. Taylor expansion proves the next root condition,
while congruence modulo $q$ preserves the full reduced Jacobian.
Construct the sequence recursively, preserving the exact parameter
coordinates and the fixed solved Jacobian. The earlier cross-level
argument proves all the stated integer congruences.
\end{proof}

This is
\lean{Komorebi.PolynomialPrescribedIntegerTower.exists_prescribed_tower_strong}.
In particular, no assumption about solvability at each later level has been
added to the initial surjectivity condition.

\begin{theorem}[Arbitrary $p$-adic parameter endpoints]
\label{thm:prescribed-padic-tower}
Let $q=p$ be prime in Theorem~\ref{thm:prescribed-integer-tower}.
For every $\tau\in t_0+p^L\Zp^k$, there are an integer tower $z_n$ and
$b\in\Zp^{m+k}$ such that
\[
 F(b)=0,\qquad\operatorname{free}(b)=\tau,\qquad
 z_n\longrightarrow b,\qquad b-z_n\in p^{L+n}\Zp^{m+k}.
\]
The initial value is $z_0=x$. At every stage and for every proposed
parameter digit $v\in\F_p^k$, the corresponding next residue fiber has
exactly $p^{m-e}$ elements. If $m=e$, then for each $\tau$ there is exactly
one such root in $B_L(x)$, independently of all choices of integer
approximants and correcting digits.
\end{theorem}

\begin{proof}
Set the first integer parameter approximation equal to the prescribed
$t_0$. For $n\ge1$, use canonical representatives of
$\tau\bmod p^{L+n}$. The initial congruence and truncation precision give
\[
 \tau-t_n\in p^{L+n}\Zp^k.
\]
Subtracting consecutive relations gives compatibility in $\Zp$.
Faithfulness of integer divisibility under $\Z\hookrightarrow\Zp$
gives compatibility in $\Z$, and the same norm bounds give $t_n\to\tau$.
Theorem~\ref{thm:prescribed-integer-tower} and completion now produce $b$
with exact polynomial vanishing and all retained precisions. Each free
coordinate has both the corresponding coordinate of $b$ and that of
$\tau$ as its limit, proving exact parameter equality.

The solved Jacobian remains surjective, so every shifted obstruction
$[\overline r+K_{z_n}v]$ vanishes. Formula~\eqref{eq:prescribed-count}
therefore gives $p^{m-e}$ choices for every parameter digit at every stage.
When $m=e$, the fixed-coordinate chart in the proof of
Theorem~\ref{thm:continuous-parameter-chart} contains every root in the
initial ball with the specified parameters. Its uniqueness proves the
endpoint assertion. This is uniqueness of the actual endpoint, not of the
integer representatives used to approach it.
\end{proof}

The endpoint theorems are
\lean{Komorebi.PolynomialPrescribedPadicTower.exists_prescribed_padic_root_with_tower}
and \lean{Komorebi.PolynomialPrescribedPadicTower.exists_unique_prescribed_padic_root}.
The separate supplied-path theorem
\lean{Komorebi.PolynomialPrescribedPadicTower.exists_root_with_prescribed_path}
accepts any compatible integer parameter path with limit $\tau$.
For $m>e$, endpoint uniqueness is false in general: the system
$F(Y_1,Y_2,T)=Y_1$ has two distinct roots $(0,0,0)$ and $(0,2,0)$ in the
same dyadic ball with identical parameter value.

\subsection{Weighted lifting for a nonunit Jacobian}

The remaining theorems in this section allow arbitrary coefficients and
centers in $R=\Zp$. They concern square solved Jacobians invertible over
$\Qp$, even when their reductions modulo $p$ are singular. The exponent
controlling an integral inverse is retained explicitly.

\begin{theorem}[Weighted nonlinear lifting]\label{thm:weighted-lifting}
Let $p$ be prime, $d,s,r\in\N$, $F\in R[X_1,\ldots,X_d]^d$, and
$a\in R^d$. Set $J=DF(a)$. Suppose
\begin{equation}\label{eq:weighted-hypotheses}
 J\in\operatorname{GL}_d(\Qp),\qquad
 p^sJ^{-1}\in M_d(R),\qquad F(a)\in p^rR^d,\qquad r>2s.
\end{equation}
Then there is exactly one root $b$ of $F$ in $a+p^{s+1}R^d$, and it
satisfies
\begin{equation}\label{eq:weighted-displacement}
 b-a\in p^{r-s}R^d.
\end{equation}
Uniqueness holds throughout the full ball of depth $s+1$.
\end{theorem}

\begin{proof}
Let $A=p^sJ^{-1}\in M_d(R)$. Inclusion into $\Qp$ and its injectivity
give the two integral identities
\[
 AJ=JA=p^sI_d.
\]
These identities are also an equivalent integral certificate for the
first two hypotheses of \eqref{eq:weighted-hypotheses}. In particular,
$A$ is injective: $Aw=0$ implies $p^sw=JAw=0$, hence $w=0$.

Precondition the polynomial system by $A$ and set $c=p^{s+1}$.
Its derivative at $a$ is $AJ=p^sI_d$. Lemma~\ref{lem:polynomial-taylor}
gives a polynomial vector $Q$ with $Q(0)=0$ and
\[
 (AF)(a+cu)=(AF)(a)+p^{2s+1}u+p^{2s+2}Q(u).
\]
Since $A$ is integral, $(AF)(a)\in p^rR^d$. The strict integer inequality
$r>2s$ permits $k=r-(2s+1)\ge0$ and a vector $C\in p^kR^d$ such that
$(AF)(a)=p^{2s+1}C$. The normalized root equation is therefore
\begin{equation}\label{eq:weighted-fixed-point}
 u=-C-pQ(u).
\end{equation}

Integral polynomial evaluation is $1$-Lipschitz for the supremum
$p$-adic metric: a common divisibility of input differences is preserved
under addition and multiplication by integral elements. Consequently the
right side of \eqref{eq:weighted-fixed-point} has contraction constant
$|p|_p<1$ on the complete space $R^d$. It has a unique fixed point $u$.

To obtain the sharper displacement, begin with $u\in R^d$. If
$u\in p^nR^d$, then $Q(0)=0$ and polynomial congruence give
$Q(u)\in p^nR^d$. For $n<k$, both $C$ and $pQ(u)$ are divisible by
$p^{n+1}$. Equation~\eqref{eq:weighted-fixed-point} therefore improves
the divisibility of $u$ by one digit. Induction gives $u\in p^kR^d$.
Thus $b=a+p^{s+1}u$ satisfies \eqref{eq:weighted-displacement}, because
$s+1+k=r-s$. Its normalized preconditioned value vanishes, and injectivity
of $A$ gives $F(b)=0$.

Every competing root $z\in a+p^{s+1}R^d$ can be written
$z=a+p^{s+1}w$. Its root equation implies the same fixed-point
equation after cancellation of the nonzero scalar $p^{2s+1}$.
Contraction uniqueness gives $w=u$, proving uniqueness on the entire
stated ball.
\end{proof}

The existence and uniqueness statements are
\lean{Komorebi.PolynomialWeightedLift.exists_weighted_root} and
\lean{Komorebi.PolynomialWeightedLift.exists_unique_weighted_root}.
The inverse hypothesis in Lean specifies an integral matrix whose image
is exactly $p^sJ^{-1}$; it does not substitute a potentially larger
determinant valuation for $s$.

\begin{example}[Strict threshold and displacement boundary]
\label{ex:weighted-boundary}
For $p=2$, $F(X)=X^2+3$, and $a=1$, one has $F(a)=4$ and $DF(a)=2$.
Thus $s=1$ and $r=2$ satisfy every displayed hypothesis except the strict
inequality. No root exists in $\Z_2$, since none exists modulo eight.
Hence $r\ge2s$ cannot replace $r>2s$ uniformly. Also,
$F(X)=p^sX+p^r$ has root $-p^{r-s}$, whose displacement is not divisible
by $p^{r-s+1}$. The guaranteed displacement cannot uniformly gain a digit.
\end{example}

\subsection{First-order coefficient and parameter motion}

\begin{theorem}[Weighted coefficient deformation]\label{thm:coefficient-motion}
Let $F,H\in R[X_1,\ldots,X_d]^d$, $F(a)=0$, and $J=DF(a)$. Suppose
$J\in\operatorname{GL}_d(\Qp)$, $p^sJ^{-1}\in M_d(R)$, and $M>2s$.
For $G=F+p^MH$, there is exactly one root $b\in a+p^{s+1}R^d$.
It satisfies $b-a\in p^{M-s}R^d$ and, in $\Qp^d$,
\begin{equation}\label{eq:coefficient-first-order}
 b-a=-p^MJ^{-1}H(a)+\varepsilon,
 \qquad\varepsilon\in p^{2M-3s}R^d.
\end{equation}
Only the original Jacobian $DF(a)$ is subject to an inverse hypothesis.
\end{theorem}

\begin{proof}
Retain the original certificate $A=p^sJ^{-1}$ and put $c=p^{s+1}$.
Taylor expansion and a polynomial difference identity yield polynomial
vectors $Q_F,Q_H$, both zero at the origin, such that
\[
 (AF)(a+cu)=p^{2s+1}u+p^{2s+2}Q_F(u),\qquad
 (AH)(a+cu)=AH(a)+cQ_H(u).
\]
Consequently
\[
 (AG)(a+cu)=p^{2s+1}\bigl(C+u+pQ(u)\bigr),
\]
where
\[
 C=p^{M-(2s+1)}AH(a),\qquad
 Q=Q_F+p^{M-s-1}Q_H,\qquad Q(0)=0.
\]
All exponents are nonnegative. The same contraction and divisibility
argument as in Theorem~\ref{thm:weighted-lifting} gives the root,
the displacement $p^{M-s}$, and full-ball uniqueness. This constructs the
root without introducing an inverse hypothesis for $DG(a)$.

For the error estimate, let $\delta=b-a$ and $D=M-s$. Taylor's theorem
for the preconditioned polynomial gives
\[
 E_F=(AF)(b)-(AF)(a)-p^s\delta\in p^{2D}R^d.
\]
Polynomial congruence gives
$E_H=(AH)(b)-(AH)(a)\in p^DR^d$.
The equations $F(a)=0$ and $G(b)=0$ imply
\[
 p^s\delta+p^MAH(a)=-E_F-p^ME_H\in p^{2D}R^d,
\]
because $M+D\ge2D$. Cancel $p^s$ in the integral domain $R$ to obtain
\[
 \delta+p^{M-s}AH(a)\in p^{2M-3s}R^d.
\]
This integral vector is $\varepsilon$. After inclusion into $\Qp$,
$p^{M-s}A=p^MJ^{-1}$ gives \eqref{eq:coefficient-first-order}.
\end{proof}

The complete theorem is
\lean{Komorebi.PolynomialRootMotion.exists_unique_root_motion}.
The error depth exceeds the displacement depth by $M-2s\ge1$. It is a
guaranteed divisibility, not an assertion of exact valuation for every
system.

\begin{theorem}[Joint coefficient and parameter deformation]
\label{thm:joint-deformation}\leavevmode\par
Let $F,H\in R[Y_1,\ldots,Y_e,T_1,\ldots,T_k]^e$ and
$a=(y_0,t_0)\in R^{e+k}$ satisfy $F(a)=0$. Define the actual blocks
\[
 J=D_YF(a)\in M_e(R),\qquad K=D_TF(a)\in M_{e\times k}(R).
\]
Suppose $J\in\operatorname{GL}_e(\Qp)$,
$p^sJ^{-1}\in M_e(R)$, and $M>2s$. For every $v\in R^k$, set
\[
 \tau=t_0+p^Mv,\qquad G=F+p^MH.
\]
There is exactly one $y\in y_0+p^{s+1}R^e$ with $G(y,\tau)=0$.
It satisfies
\begin{equation}\label{eq:joint-first-order}
 \begin{aligned}
  y-y_0&\in p^{M-s}R^e,\\
  y-y_0&=-p^MJ^{-1}\bigl(H(a)+Kv\bigr)+\varepsilon,
  \qquad\varepsilon\in p^{2M-3s}R^e.
 \end{aligned}
\end{equation}
The equality involving $J^{-1}$ is in $\Qp^e$, with an integral error
vector. The parameter equality is exact.
\end{theorem}

\begin{proof}
Use the augmented polynomial systems
\[
 \mathcal F(Y,T)=(F(Y,T),T-t_0),\qquad
 \mathcal H(Y,T)=(H(Y,T),-v).
\]
Then $\mathcal F(a)=0$, and its derivative is
\[
 B=D\mathcal F(a)=\begin{pmatrix}J&K\\0&I_k\end{pmatrix}.
\]
For $A=p^sJ^{-1}\in M_e(R)$ and $c=p^s$, define
\[
 C=\begin{pmatrix}A&-AK\\0&cI_k\end{pmatrix}.
\]
The identities $AJ=JA=cI_e$ give $CB=BC=cI_{e+k}$. In the second
product, the upper right block is $-JAK+cK=0$; in the first, it is
$AK-AK=0$. Thus augmentation preserves the original inverse bound $s$.

Theorem~\ref{thm:coefficient-motion} gives a unique root $b$ of
$\mathcal F+p^M\mathcal H$ in $a+p^{s+1}R^{e+k}$, with displacement
depth $M-s$. Its last equations say exactly
$b_T=t_0+p^Mv=\tau$. Its first equations are $G(b)=0$.
The integral first-order expression uses
\[
 C\binom{H(a)}{-v}
 =\binom{A(H(a)+Kv)}{-p^sv}.
\]
Projection onto the solved coordinates gives
\[
 b_Y-y_0+p^{M-s}A(H(a)+Kv)\in p^{2M-3s}R^e.
\]
The same inclusion identity as before converts this to
\eqref{eq:joint-first-order}.

Finally, if $z\in y_0+p^{s+1}R^e$ is any competing solved root with
parameters $\tau$, then $(z,\tau)$ belongs to the full augmented ball.
Indeed, $M>2s$ implies $M\ge s+1$. Full-point uniqueness therefore gives
$(z,\tau)=b$. This comparison proves uniqueness on exactly the solved
ball stated in the theorem.
\end{proof}

The exact solved-coordinate statement is
\lean{Komorebi.PolynomialJointDeformation.exists_joint_deformation},
with its unique-existence interface
\lean{Komorebi.PolynomialJointDeformation.exists_unique_joint_deformation}.
The triangular inverse identities hold over arbitrary commutative rings,
so they apply in particular over the $p$-adic integers.

\begin{example}[First-order sign and error precision]\label{ex:motion-boundary}
For $p=2$, $s=0$, $M=2$, $F(Y,T)=Y+T$, $H=1$, $a=0$, and $v=1$,
the prescribed parameter is $\tau=4$ and the solved root is $y=-8$.
Here $J=K=1$, and the forcing term is $H(a)+Kv=2$. Omitting the parameter
term would predict $-4$, and reversing its sign would predict $0$.
The corresponding errors $-4$ and $-8$ are both outside
$2^{2M}\Z_2=16\Z_2$. Thus either change would violate the stated error
bound.

The error exponent cannot uniformly be increased. For
$F(Y,T)=Y+Y^2+T$, $H=0$, $a=0$, $p=2$, $s=0$, $M=1$, and $v=1$,
there is an even root $y$ of $y+y^2+2=0$. Its equation forces
$y\equiv2\pmod4$ and $y+2=-y^2$, so
\[
 4\mid y+2,\qquad 8\nmid y+2.
\]
Thus the error depth $2M-3s=2$ is attained even for parameter-only motion.
The same quadratic with a constant coefficient perturbation gives the
coefficient-only boundary. These examples refute a uniform extra error
digit; they do not assert optimality for every individual choice of
$p,s,M$.
\end{example}

The combined results establish finite singular obstructions, smooth
compatible towers, full continuous root charts after a suitable coordinate
choice, and quantitative deformation for a rationally invertible solved
Jacobian. Rationally rank-deficient persistence remains a separate problem.
The analytic solution-map theorem in the next section supplies the
regularity conclusion for these same derivative blocks.
